\TieudeT{PHONG TỎA VẬT LÍ}
\subsubsection*{PHẦN I. Thí sinh trả lời câu hỏi từ câu 1 đến câu 18. Mỗi câu hỏi thí sinh chỉ chọn một phương án}
\setcounter{ex}{0}
\begin{ex}%book
Một vật rơi tự do sau thời gian $t$ đi được quãng đường $s$, gia tốc rơi tự do được tính theo công thức nào sau đây?
\choice
{$g=\dfrac{2 s}{t}$}
{$g=\dfrac{s}{2 t^2}$}
{\True $g=\dfrac{2 s}{t^2}$}
{$g=\dfrac{s}{t^2}$}

\loigiai{
}
\haidong{7} \end{ex}
\begin{ex}%book
Sự rơi tự do là chuyển động rơi
\choice
{không chịu tác dụng của bất kì lực nào}
{chịu tác dụng của lực cản không khí}
{không chịu tác dụng của lực cản không khí}
{\True chỉ chịu tác dụng của trọng lực}

\loigiai{
}
\haidong{7} \end{ex}
\begin{ex}%book
Đặc điểm nào sau đây \textbf{không} phải là của chuyển động rơi tự do?
\choice
{\True Chuyển động thẳng chậm dần đều}
{Chuyển động theo phương thẳng đứng, chiều từ trên xuống}
{Chuyển động nhanh dần đều}
{Tại một vị trí xác định và ở gần mặt đất, mọi vật rơi tự do là như nhau}
\loigiai{
}
\haidong{7} \end{ex}
\begin{ex}%book
Phát biểu nào sau đây là \textbf{sai}?  Rơi tự do
\choice
{là chuyển động nhanh dần đều}
{\True là chuyển động thẳng đều}
{phương thẳng đứng}
{chiều từ trên xuống}

\loigiai{
}
\haidong{7} \end{ex}
\begin{ex}%book
Phát biểu nào sau đây là \textbf{sai}? Chuyển động rơi tự do có
\choice
{phương thẳng đứng}
{chiều từ trên xuống}
{gia tốc không đổi}
{\True vận tốc không đổi}
\loigiai{
}
\haidong{7} \end{ex}
\begin{ex}%book %[1P1B7-1][Loai1] 
Chuyển động của vật nào dưới đây sẽ được coi gần đúng nhất là rơi tự do nếu được thả rơi?
\choice
{Một cái lá cây}
{Một sợi chỉ}
{Một chiếc khăn tay}
{\True Một viên bi sắt}
\loigiai{
\begin{itemize}
\item Rơi tự do là chỉ do tác dụng của trọng lực.
\item Trong các trường hợp này thì cục phấn khi thả rơi là trường hợp có rơi tự do.
\end{itemize}
}
\haidong{7} \end{ex}
\begin{ex}%book %[1P1B7-1][Loai1] 
Trường hợp nào sau đây \textbf{không} thể coi chuyển động của vật là sự rơi tự do?
\choice
{Một chiếc lông chim rơi trong ống hút chân không}
{Một hòn đá rơi từ trên mái nhà xuống với vận tốc đầu bằng không}
{\True Một người nhảy dù khi đã căng dù}
{Một quả táo rụng từ trên cây xuống}
\loigiai{
Một người nhảy dù khi đã căng dù không thể coi là sự rơi tự do vì trường hợp này lực cản của không khí đáng kể so với trọng lực tác dụng lên người và dù.
}
\haidong{7} \end{ex}

\begin{ex}%book
Hai vật có khối lượng $m_1>m_2$ rơi tự do tại cùng một địa điểm và cùng một độ cao thì vận tốc chạm đất 
\choice
{$v_1>v_2$}
{\True $v_1=v_2$}
{$v_1<v_2$}
{Không có cơ sở kết luận}
\loigiai{
}
\haidong{7} \end{ex}
\begin{ex}%book
Một vật rơi tự do từ độ cao $h_1$ đến mặt đất trong thời gian $t_1$. Một vật khác rơi tự do từ độ cao $h_2$ đến mặt đất trong thời gian $t_2$. Hệ thức nào sau đây là đúng?
\choice
{$\dfrac{h_1}{h_2}=\dfrac{t_1}{t_2}$}
{$\dfrac{h_1}{h_2}=\dfrac{t_2}{t_1}$}
{$\dfrac{h_1}{h_2}=\left(\dfrac{t_2}{t_1}\right)^2$}
{\True $\dfrac{h_1}{h_2}=\left(\dfrac{t_1}{t_2}\right)^2$}

\loigiai{
Ta có $\left\{\begin{array}{l}h_1=\dfrac{1}{2} g t_1^2 \\ h_2=\dfrac{1}{2} g t_2^2\end{array} \Rightarrow \dfrac{h_1}{h_2}=\left(\dfrac{t_1}{t_2}\right)^2\right.$}
\haidong{7} %<MyLT1>
\end{ex}

\begin{ex}%book
Trong sự rơi tự do thì vận tốc của vật rơi 
\choice
{\True tỉ lệ thuận với căn bậc hai của đoạn đường rơi}
{tỉ lệ nghịch với đoạn đường rơi}
{tỉ lệ thuận với đoạn đường rơi}
{tỉ lệ thuận với bình phương của đoạn đường rơi}
\loigiai{
}
\haidong{7} \end{ex}
\begin{ex}%book%[1P1B7-2][Loai1]
Một vật được thả từ trên máy bay ở độ cao $80\ m$. Bỏ qua sức cản không khí. Lấy $g = 10\ m/s^2$. Thời gian rơi là	
\choice
{$4,04\ s$}
{$8,00\ s$}
{\True $4,00\ s$}
{$2,86\ s$}	
\haidong{2}
\loigiai{
\begin{itemize}
\item Ta có: $S =\dfrac{gt^2}{2}=5t^2 $.
\item Khi vật chạm đất: $h=S = 80\ m \Rightarrow 5t^2=80 \Rightarrow t=4\ s$.
\end{itemize}	
}
\haidong{10} \end{ex}
\begin{ex}%book
Một hòn đá rơi từ miệng một giếng cạn đến đáy giếng mất $3\ s$. Lấy $g=9,8\ m/s^2$. Độ sâu của giếng ước lượng là
\choice
{$90\ m$}
{$45\ m$}
{$30,4\ m$}
{\True $44,1\ m$}	
\haidong{2}
\loigiai{
Độ sâu của giếng là $h=s=\dfrac{1}{2} g \cdot t^2=44,1\ m$.		
}
\haidong{10} \end{ex}

\begin{ex}%book%[1P1K7-2][Loai1]
Một người thả vật rơi tự do, vật chạm đất có $v = 30\ m/s$. Lấy $g = 10\ m/s^2$. Độ cao của vật sau khi đi được $2\ s$ là
\choice
{$30\ m$}
{\True $25\ m$}
{$20\ m$}
{$15\ m$}
\haidong{4}
\loigiai{
\begin{itemize}
\item Chọn gốc tọa độ tại vị trí vật rơi, chiều dương hướng xuống.
\item Ta có: $v_t=gt=10t;\ y=\dfrac{gt^2}{2}=5t^2$.
\item Vận tốc chạm đất là $v = 30\ m/s \Rightarrow$ thời gian vật rơi là: $t=\dfrac{v_t}{10}=3\ s$.  
\item Độ cao ban đầu của vật là: $y_3=5.3^2=45\ m$.
\item Quãng đường vật đi được sau 2 s là: $y_2=5.2^2=20\ m$.
\item Độ cao của vật sau khi đi được 2 s là: $h=y_3-y_2=45-20=25\ m$.
\end{itemize}	
}
\haidong{10} \end{ex}
\begin{ex}%book %[1P1-7.2K][Loai1] 
Hai hòn đá được thả rơi tự do từ cùng một độ cao nhưng sớm, muộn hơn nhau $1\ s$. Khi hòn đá thả trước chạm đất thì hòn đá thả sau còn cách mặt đất $35\ m$. Lấy $g = 10\  m/s^2$. Chiều cao của hai hòn đá lúc ban đầu là
\choice
{$75\ m$}
{\True $80\ m$}
{$85\ m$}
{$90\ m$}
\haidong{6}
\loigiai{		
\begin{itemize}
\item Gọi h là chiều cao của hai hòn đá; t là thời gian rơi của hòn đá rơi trước.
\item Với hòn đá rơi trước: $h_1 = \dfrac{gt^2}{2}$.
\item Hòn đá rơi sau: $h_2 = \dfrac{g(t - 1)^2}{2} \Rightarrow h - h_1 = 35 = \dfrac{g}{2}(t^2 - (t - 1)^2) \Rightarrow t = 4\ s$ hay $h = 80$ m.
\end{itemize}	
}
\haidong{10} \end{ex}
\begin{ex}%book %[1P1K7-3][Loai1] %BAI 4+5 - HOC MAI
Một vật rơi từ độ cao $h$, trong giây cuối cùng vật rơi được quãng đường $25\ m$. Lấy $g = 10\  m/s^2$. Thời gian vật rơi tự do từ độ cao $h$ là
\choice
{$1\ s$}
{$2\ s$}
{\True $3\ s$}
{$4\ s$}
\haidong{6}
\loigiai{  
\begin{itemize}
\item Khi vật rơi trong thời gian t: $h_1 = \dfrac{1}{2}gt^2 $
\item Với $h_1 = h$
\item Khi vật rơi trong thời gian $t-1$: $h_2 = \dfrac{1}{2}g(t-1)^2 $
\item Ta có: $\Delta h = h_1 - h_2 = \dfrac{1}{2}g\left( {t^2 - \left( {t - 1} \right)^2} \right) = 25 \Rightarrow t = 3\ s.$
\end{itemize}	}
\haidong{10} \end{ex}
\begin{ex}%book%[1P1K7-2][Loai1]
Để biết độ sâu của một cái hang, những người thám hiểm thả một hòn đá từ miệng hang và đo thời gian từ lúc thả đến lúc nghe thấy tiếng vọng của hòn đá khi chạm đất. Giả sử người ta đo được thời gian là $13,66\ s$. Tính độ sâu của hang. Lấy gia tốc trọng trường $g = 10\ m/s^2$ và vận tốc âm trong không khí là $340\ m/s$. Coi âm truyền theo một phương nào đó là thẳng đều.
\choice
{\True $680\ m$}
{$450\ m$}
{$380\ m$}
{$820\ m$}
\haidong{9}
\loigiai{
\begin{itemize}
\item Thời gian rơi tự do $t_1$ của hòn đá: $h=\dfrac{1}{2}gt_1^2\Rightarrow t_1=\sqrt{{\dfrac{2h}g}}=\sqrt{{\dfrac{h}{5}}}$.
\item Thời gian truyền âm $t_2$ kể từ khi hòn đá chạm đất đến khi nghe thấy tiếng vọng: $t_2=\dfrac{h}{v}=\dfrac{h}{340}$.
\item Theo đề ta có: $t_1+t_2=13,66\Leftrightarrow \dfrac{h}{340}+\sqrt{{\dfrac{h}{5}}}=13,66\Rightarrow h\approx 680\ m$.
\end{itemize}
}
\haidong{10} \end{ex}

\begin{ex} %[1P1-11.2Y][Loai1] 
\immini[thm]{Dùng một đồng hồ đo thời gian có độ chia nhỏ nhất $0,001\ s$ để đo $n$ lần thời gian rơi tự do của một vật bắt đầu từ điểm $A$ ($v_A = 0$) đến điểm $B$, kết quả cho trong bảng. Thời gian rơi trung bình là
\motcot
{\True $\overline{t}=0,404\ s$}
{$\overline{t}=0,410\ s$}
{$\overline{t}=0,398\ s$}
{$\overline{t}=0,402\ s$}}{$\begin{array} { | c | c |} \hline n &  t  \\ \hline 1 & { 0,398 } \\ \hline 2 & { 0,399 } \\ \hline 3 & { 0,408 } \\ \hline 4 & { 0,406 } \\ \hline 5 & { 0,405 } \\ \hline 6 & { 0,402 } \\ \hline 7 & { 0,410 } \\ \hline \end{array}$}
\loigiai{
Thời gian rơi trung bình: $\overline t=\dfrac{0,398+0,399+0,408+0,406+0,405+0,402+0,410}{7}=0,404\ s$
}
\haidong{10} \end{ex}
\begin{ex}
Trong giờ thực hành, một học sinh đo thời gian rơi của một vật từ cùng một độ cao. Kết quả $5$ lần đo được cho ở bảng sau
\begin{center}
\begin{tabular}{|l|l|l|l|l|l|}
\hline
Lần đo & 1 & 2 & 3 & 4 & 5 \\
\hline
Thời gian (s) & 2,48 & 2,49 & 2,46 & 2,50 & 2,47 \\
\hline
\end{tabular}
\end{center}
Cho thang chia nhỏ nhất của đồng hồ là $0,02\ s$. Lấy sai số dụng cụ bằng nửa độ chia nhỏ nhất. Sai số tuyệt đối của phép đo là
\choice
{$0,01\ s$}
{\True $0,02\ s$}
{$0,04\ s$}
{$0,03\ s$}
\haidong{4}

\loigiai{
}
\haidong{10} \end{ex}


















\subsubsection*{PHẦN 2. Thí sinh trả lời từ câu 1 đến câu 4. Trong mỗi ý a), b), c), d) ở mỗi câu, thí sinh chọn đúng hoặc sai.}
\setcounter{ex}{0}
\begin{ex}%Tailieuchuan
Một vật được thả rơi tự do từ độ $h$ so với mặt đất. Gia tốc rơi tự do là $g$.
\choiceTF[t]
{\True Vật rơi nhanh dần đều} 
{\True Thời gian vật rơi hết độ cao $h$ là $t=\sqrt{\dfrac{2h}{g}}$} 
{Tốc độ vật chạm đất là $v=\sqrt{gh}$} 
{Quãng đường vật rơi được trong giây thứ $n$ là $\Delta s_{n}=(2 n-1)\ g$} 
\loigiai{ 
\begin{itemchoice}
\itemch
\itemch
\itemch
\itemch
\end{itemchoice}
}
\end{ex}

\begin{ex}%Tailieuchuan
Thời gian rơi của một vật được thả rơi tự do tới mặt đất là $8\ s$. Lấy $g=10\ m/s^2$.
\choiceTF[t]
{Trong quá trình rơi vận tốc của vật không thay đổi} 
{\True Vận tốc của vật ngay trước khi chạm đất là $80\ m/s$} 
{Độ cao nơi thả vật so với mặt đất là $640\ m$} 
{Khi cách mặt đất $75\ m$ vật có vận tốc $75 \ m/s$} 
\loigiai{
\begin{itemchoice}
\itemch
\itemch
\itemch
\itemch
\end{itemchoice}
}
\end{ex}
\begin{ex}%Tailieuchuan
Một vật rơi tự do tại một địa điểm có độ cao $45\ m$ xuống mặt đất, lấy $g=10 \ m/s^2$.
\choiceTF[t]
{\True Phương rơi của vật là phương thẳng đứng} 
{Thời gian vật chạm đất là $2,5\ s$} 
{\True Quãng đường vật rơi được sau $2\ s$ là $20\ m$} 
{\True Thời gian vật rơi hết $13,75 \ m$ cuối cùng là $0,5\ s$} 
\loigiai{ 
\begin{itemchoice}
\itemch
\itemch
\itemch
\itemch
\end{itemchoice}
}
\end{ex}
\begin{ex}%Tailieuchuan
Một vật rơi tự do, trong $2\ s$ cuối cùng trước khi chạm đất đi được quãng đường $180\ m$. Lấy $g=10\ m/s^2$.
\choiceTF[t]
{Trong quá trình rơi gia tốc của vật tăng dần} 
{Thời gian vật rơi từ vị trí thả vật là $9\ s$} 
{\True Vật được thả rơi tự do từ nơi có độ cao $500\ m$ so với mặt đất} 
{\True Thời gian vật rơi trong $255\ m$ cuối cùng là $3\ s$} 
\loigiai{
\begin{itemchoice}
\itemch
\itemch
\itemch
\itemch
\end{itemchoice}
}
\end{ex}












\subsubsection*{PHẦN 3. Thí sinh trả lời từ câu 1 đến câu 6.}
\setcounter{ex}{0} 
\begin{ex}%Tailieuchuan
Một vật rơi tự do từ độ cao $h$ xuống đất. Lấy $g=10\ m/s^2$. Quãng đường vật rơi được trong $5\ s$ đầu tiên là bao nhiêu mét?
\shortans[oly]{125}
\loigiai{
}
\end{ex}
\begin{ex}
Một vật được thả rơi tự do từ độ cao $h$ so với mặt đất, tại nơi có gia tốc trọng trường là $g=10 \ m/s^2$.
Tốc độ của vật khi chạm đất là $120 \ m/s$. Tìm $h$.
\shortans[oly]{720}
\loigiai{
Khi vật chạm đất thì độ cao nơi thả vật bằng quãng đường vật rơi.

Vì tốc độ của vật khi chạm đất bằng $120 \ m/s$ nên độ cao nơi thả vật
$$
h=s=\dfrac{v^2}{2 \ g}=\dfrac{120^2}{2.10}=720 \ m
$$
}
\end{ex}

\begin{ex}%Tailieuchuan
Một vật rơi tự do từ độ cao $80\ m$ xuống đất. Lấy $g=10\ m/s^2$. Quãng đường vật rơi được trong $0,5\ s$ cuối cùng là bao nhiêu mét (làm tròn kết quả đến chữ số hàng phần mười)?
\shortans[oly]{18,8}
\loigiai{
}
\end{ex}
\begin{ex}
Một viên bi sắt được thả rơi từ trên cao xuống mặt đất trong thời gian $5,5 \ s$, tại nơi có gia tốc trọng trường $g=10 \ m/s^2$. Bỏ qua sức cản của không khí. Thời gian rơi của viên bi kể từ khi thả cho đến khi nó còn cách mặt đất $15\ m$ là bao nhiêu giây (làm tròn kết quả đến chữ số hàng phần trăm)? \shortans[oly]{5,22}
\loigiai{

Độ cao nơi thả vật: $h=\dfrac{1}{2}$ g. $t^2=\dfrac{1}{2} \cdot 10 \cdot 6^2=180 \ m$
Khi vật rơi cách mặt đất 15 m thì quãng đường mà vật rơi được là: $s^{\prime}=h-15=180-15=165 \ m$.

Thời gian khi vật rơi cách mặt đất 15 m là:

$$
s^{\prime}=\dfrac{1}{2} gt^{, 2} \Rightarrow t^{\prime}=\sqrt{\dfrac{2. s^{\prime}}{g}}=\sqrt{\dfrac{2.165}{10}}=5,74 \ s.
$$

}
\end{ex}

\begin{ex}%Tailieuchuan
Từ một đỉnh tháp cao $50\ m$, người ta buông vật thứ nhất. Sau $1\ s$ thì người ta lại buông vật thứ hai ở tầng thấp hơn đỉnh tháp $26,5\ m$. Lấy $g=10\ m/s^2$. Bỏ qua sức cản không khí. Thời gian kể từ lúc vật thứ nhất rơi đến khi hai vật gặp nhau là bao nhiêu giây?
\shortans[oly]{3,15}
\loigiai{
}
\end{ex}

\begin{ex}%Tailieuchuan
Hai viên bi rơi tự do từ cùng một độ cao cách nhau một khoảng thời gian $1,5\ s$. Lấy $g=10\ m/s^2$. Khi chưa có viên bi nào chạm đất thì khoảng cách giữa hai viên bi sau khi viên bi thứ nhất rơi được $3,5\ s$ là bao nhiêu mét (làm tròn kết quả đến chữ số hàng phần mười)?
\shortans[oly]{41,3}
\loigiai{}
\end{ex}

\newpage 

\TieudeT{PHONG TỎA VẬT LÍ}
\subsubsection*{PHẦN I. Thí sinh trả lời câu hỏi từ câu 1 đến câu 18. Mỗi câu hỏi thí sinh chỉ chọn một phương án}
\setcounter{ex}{0}
\begin{ex}%book
Một vật rơi tự do từ độ cao h xuống đất. Tính thời gian rơi của vật đến khi chạm đất được tình theo công thức nào sau đây?
\choice
{$t=2 gh$}
{\True $t=\sqrt{\dfrac{2 h}{g}}$}
{$t=\sqrt{gh}$}
{$t=\sqrt{2 gh}$}

%\loigiai{
%}
\haidong{7} \end{ex}
\begin{ex}%book
Đặc điểm nào sau đây phù hợp với sự rơi tự do?
\choice
{Chuyển động thẳng đều}
{Lực cản của không khí lớn}
{\True Vận tốc được tính theo công thức $v = gt$}
{Vận tốc giảm dần theo thời gian}
%\loigiai{
%}
\haidong{7} \end{ex}
\begin{ex}%book
Chuyển động rơi tự do là một chuyển động thẳng
\choice
{chậm dần đều}
{đều}
{\True nhanh dần đều}
{biến đổi không đều}

%\loigiai{
%}
\haidong{7} \end{ex}
\begin{ex}%book
Phát biểu nào sau đây là \textbf{sai}?  Chuyển động rơi tự do
\choice
{có phương thẳng đứng}
{có chiều từ trên xuống}
{là chuyển động thẳng, nhanh dần đều}
{\True vào lúc $t=0$ thì $v \neq 0$}

%\loigiai{
%}
\haidong{7} \end{ex}
\begin{ex}%book
Phát biểu nào sau đây là \textbf{sai}?
\choice
{Khi rơi tự do tốc độ của vật tăng dần}
{Vật được coi là rơi tự do khi lực cản không khí rất nhỏ so với trọng lực}
{\True Vận động viên nhảy dù từ máy bay xuống mặt đất sẽ rơi tự do}
{Rơi tự do có quỹ đạo là đường thẳng}
%\loigiai{
%}
\haidong{7} \end{ex}
\begin{ex}%book %[1P1B7-1][Loai1] 
Chuyển động nào dưới đây có thể coi như là chuyển động rơi tự do?
\choice
{Chuyển động của một hòn sỏi được ném lên cao}
{Chuyển động của một hòn sỏi được ném theo phương nằm ngang}
{Chuyển động của một hòn sỏi được ném phương xiên góc}
{\True Chuyển động của một hòn sỏi được thả rơi xuống}
%\loigiai{
%Rơi tự do là sự rơi do tác dụng của trọng lực.
%}
\haidong{7} \end{ex}
\begin{ex}%book
Gia tốc rơi tự do phụ thuộc vào những yếu tố nào?
\choice
{Khối lượng và kích thước vật rơi}
{\True Độ cao và vĩ độ địa lí}
{Vận tốc đầu và thời gian rơi}
{Áp suất và nhiệt độ môi trường}
%\loigiai{
%}
\haidong{7} %<MyLT1>
\end{ex}
\begin{ex}%book
Một vật rơi trong không khí nhanh chậm khác nhau, nguyên nhân nào sau đây quyết định điều đó?
\choice
{Do các vật làm bằng chất liệu khác nhau}
{Do các vật to nhỏ khác nhau}
{\True Do lực cản của không khí lên các vật}
{Do các vật nặng nhẹ khác nhau}
%\loigiai{
%}
\haidong{7} \end{ex}
\begin{ex}%book
Thả hai vật rơi tự do không vận tốc đầu, đồng thời từ hai cao $h_1$ và $h_2$ so với mặt đất tại cùng một nơi trên Trái Đất. Nếu thời gian từ lúc rơi đến lúc chạm đất của vật thứ nhất bằng $2$ lần vật thứ $2$ thì
\choice
{$h_1=h_2$}
{\True $h_1=4\ h_2$}
{$h_1=3\ h_2$}
{$h_1=2,5\ h_2$}

%\loigiai{
%}
\haidong{7} \end{ex}
\begin{ex}%book
Một vật rơi tự do, nếu thời gian rơi tăng lên $2$ lần thì quãng đường đi của vật sẽ
\choice
{tăng $2$ lần}
{\True tăng $4$ lần}
{tăng $1,4$ lần}
{tăng $2,8$ lần}

\loigiai{
Ta có $s=\dfrac{1}{2} g t^2$

Thấy $s \sim t^2$ nên t tăng 4 lần thì s tăng 4 lần.

}
\haidong{7} \end{ex}

\begin{ex}%book
Trong chuyển động rơi tự do, đồ thị phụ thuộc của độ dịch chuyển $(d)$ theo thời gian rơi $(t)$ có dạng là
\choice
{đường thẳng song song với trục $Ot$}
{đường thẳng qua gốc tọa độ}
{đường thẳng không qua gốc tọa độ}
{\True một nhánh của đường parabol}
\loigiai{
Ta có độ dịch chuyển $d=s=\dfrac{1}{2} g t^2$
$\Rightarrow$ đồ thị của $(d)$ theo $(t)$ là một nhánh của đường Parabol $(t \geq 0)$.
}
\haidong{7} \end{ex}
\begin{ex}%book %[1P1B7-2][Loai1]
Một vật nặng rơi tự do từ độ cao $20\ m$ xuống đất. Lấy $g = 10\  m/s^2$. Tốc độ của vật khi chạm đất là
\choice
{$10\sqrt 2\   m/s$}
{\True $20\  m/s$}
{$400\  m/s$}
{$4\  m/s$}
\haidong{2}
\loigiai{		
\begin{itemize}
\item Từ công thức liên hệ giữa độ dời, vận tốc và gia tốc: $v_t^2 - v_0^2 = 2as$.
\item Ta áp dụng vào sự rơi tự do: $\Delta y = h;\ a = g;\ v_0 = 0$.
\item Ta có: $v_t = \sqrt {2gh} = \sqrt {2.10.20} = 20\  m/s^2$.
\end{itemize}	
}
\haidong{10} \end{ex}
\begin{ex}%book
Lấy gia tốc rơi tự do là $g=10\ m/s^2$. Biết thời gian rơi tự do của vật là $2$ giây. Độ dịch chuyển của một vật khi đó là 	
\choice
{$25\ m$}
{$10\ m$}
{$15\ m$}
{\True $20\ m$}
\haidong{2}

\loigiai{
}
\haidong{10} \end{ex}
\begin{ex}%book
Thời gian rơi tự do của một vật là $6\ s$. Lấy $g=10\ m/s^2$. Tốc độ lúc chạm đất của vật là
\choice
{$40\ m/s$}
{$50\ m/s$}
{\True $60\ m/s$}
{$80\ m/s$}
\haidong{2}
\loigiai{
Tốc độ của vật ngay trước khi chạm đất là $v=g. t=60\ m/s$.

}
\haidong{10} \end{ex}

\begin{ex}%book %[1P1K7-3][Loai1] %BAI 4+5 - HOC MAI
Thả một hòn sỏi từ trên gác cao xuống đất. Trong giây cuối cùng hòn sỏi rơi được quãng đường $15\ m$. Lấy $g = 10\  m/s^2$. Độ cao của điểm từ đó bắt đầu thả hòn sỏi là
\choice
{$15\ m$}
{$10\ m$}
{\True $20\ m$}
{$45\ m$}
\haidong{6}
\loigiai{

\begin{itemize}
\item Quãng đường vật đi được trong thời gian t: $S = \dfrac{1}{2}gt^2$
\item Quãng đường vật đi được trong thời gian $(t - 1)$: $S' = \dfrac{1}{2}g(t - 1)^2$ 
\item Quãng đường vật đi được trong giây cuối cùng: \begin{align*}\Delta S = S - S' = \dfrac{1}{2}gt^2 - \dfrac{1}{2}g{\left( {t - 1} \right)^2} = gt - \dfrac{1}{2}g = 15\ m \Rightarrow t = 2\ s \Rightarrow S = \dfrac{1}{2}gt^2 = 20\ m.\end{align*}
\end{itemize}	}
\haidong{10} \end{ex}
\begin{ex}%book
Một vật được thả rơi tự do từ độ cao $h$. Trong nửa đoạn đường đầu, vật rơi trong thời gian $t_1$;\ trong nửa đoạn đường sau, vật rơi trong thời gian $t_2$. Hệ thức nào sau đây là đúng?
\choice
{$\dfrac{t_2}{t_1}=\sqrt{2}$}
{\True $\dfrac{t_2}{t_1}=\sqrt{2}-1$}
{$\dfrac{t_2}{t_1}=1$}
{$\dfrac{t_2}{t_1}=2$}
\haidong{6}
\loigiai{
\immini[thm]{Vật rơi trong nửa đoạn đường đầu với thời gian $t_1=\sqrt{\dfrac{2 s_1}{g}}$

Vật rơi trong hết đoạn đường với thời gian $t=\sqrt{\dfrac{2 s}{g}}=\sqrt{\dfrac{2 \cdot 2 s_1}{g}}=t_1 \cdot \sqrt{2}$

$\rightarrow$ Vật rơi trong nửa đoạn đường sau với thời gian là $t_2=t-t_1=t_1(\sqrt{2}-1)$

$\Rightarrow \dfrac{t_2}{t_1}=\sqrt{2}-1$}{\includegraphics[width=2cm]{images/8_1}}

}
\haidong{10} \end{ex}
\begin{ex}
Trong một bài thực hành, gia tốc rơi tự do được tính theo công thức $g=\dfrac{2h}{t^2}$. Sai số tỉ đối của phép đo trên tính theo công thức nào?
\choice
{$\dfrac{\Delta g}{\overline{g}}=\dfrac{\Delta h}{\overline{h}}+\dfrac{\Delta t}{\overline{t}}$}
{\True $\dfrac{\Delta g}{\overline{g}}=\dfrac{\Delta h}{\overline{h}}+2 \dfrac{\Delta t}{\overline{t}}$}
{$\dfrac{\Delta g}{\overline{g}}=\dfrac{\Delta h}{\overline h}\pm 2 \dfrac{\Delta t}{\overline{t}}$}
{$\dfrac{\Delta g}{\overline{g}}=\dfrac{\Delta h}{\overline{h}}-2 \dfrac{\Delta t}{\overline{t}}$}

\loigiai{
}
\haidong{10} \end{ex}

\begin{ex} %[1P1-11.2B][Loai1] 
\immini[thm]{Dùng một đồng hồ đo thời gian có độ chia nhỏ nhất 0,001 s để đo n lần thời gian rơi tự do của một vật bắt đầu từ điểm $A$ ($v_A = 0$) đến điểm $B$, kết quả cho trong bảng. Sai số ngẫu nhiên của phép đo là
\motcot
{\True $\overline{\Delta t}=0,004\ s$}
{$\overline{\Delta t}=0,410\ s$}
{$\overline{\Delta t}=0,005\ s$}
{$\overline{\Delta t}=0,003\ s$}}
{$\begin{array} { | c | c | } \hline \text { n } & { t } \\ \hline 1 & { 0,398 } \\ \hline 2 & { 0,399 } \\ \hline 3 & { 0,408 } \\ \hline 4 & { 0,406 } \\ \hline 5 & { 0,405 } \\ \hline 6 & { 0,402 } \\ \hline
7 & { 0,410 } \\ \hline \end{array}$}
\haidong{4}
\loigiai{
\begin{itemize}
\item Giá trị trung bình của thời gian: $\overline t=\dfrac{0,398+0,399+0,408+0,406+0,405+0,402+0,410}{7}=0,404\ s$.
\item Sai số ngẫu nhiên của phép đo: $\overline{\Delta t}=\dfrac{0,006+0,005+0,004+0,002+0,004+0,002+0,006}{7}=0,004\ s$.
\end{itemize}
}
\haidong{10} \end{ex}










\subsubsection*{PHẦN 2. Thí sinh trả lời từ câu 1 đến câu 4. Trong mỗi ý a), b), c), d) ở mỗi câu, thí sinh chọn đúng hoặc sai.}
\setcounter{ex}{0}
\begin{ex}
\immini[thm]{Cho thiết bị thí nghiệm đo gia tốc trọng trường như hình vẽ.
\choiceTF[t]
{Số $1$ là công tắc cổng quang}
{Đo gia tốc trọng trường là phép đo trực tiếp}
{Để đo quãng thời gian qua cổng $A$ và $B$ thì ta chỉnh đồng hồ về MODE A}
{\True Số $3$ là đồng hồ đo thời gian}
} 
{\includegraphics[width=4cm]{images/2025_06_04_7bf78a016ca571c65746g-173}
} 

\loigiai{
\begin{itemchoice}
\itemch
\itemch
\itemch
\itemch
\end{itemchoice}}
\end{ex}



\begin{ex}%book
Một người thả rơi tự do một hòn bi từ trên cao xuống đất và đo được thời gian rơi là $3\ s$. Bỏ qua sức cản không khí. Lấy $g=10 \ m / s^{2}$.
\choiceTF[t]
{\True Chuyển động của hòn bi rơi là chuyển động thẳng nhanh dần đều}
{\True Độ cao của vị trí thả hòn bi so với mặt đất là $45\ m$}
{Quãng đường hòn bi rơi được trong $0,5 \ s$ cuối trước khi chạm đât là $1,25 \ m$}
{\True Vận tốc của hòn bi lúc chạm đất là $30 \ m / s$}			
\loigiai{\begin{itemchoice}
\itemch Đúng.
\itemch Đúng. Độ cao của nơi thả hòn bi so với mặt đất là $h=\dfrac{1}{2} g t^{2}=\dfrac{1}{2} . 10 . 3^{2}=45(\ m)$
\itemch Sai. Quãng đường hòn bi rơi được trong $0,5 \ s$ cuối trước khi chạm đất $=$ cả quãng đường - quãng đường vật rơi trong 2,5 giây đầu: $h_{t=0,5 \ s}=h-h_{t=2,5 \ s}=45-\dfrac{1}{2} . 10 . 2,5^{2}=13,75\ (m)$
\itemch Đúng. Vận tốc của hòn bi lúc chạm đất là $v=g . t=10.3=30\ (m / s)$
\end{itemchoice}
}

\haidong{10} \end{ex}
\begin{ex}
Một học sinh thả một vật rơi tự do từ đỉnh tháp có độ cao $125\ m$ so với mặt đất, tại nơi có gia tốc trọng trường $g=10 \ m/s^2$. Chọn mốc thời gian $t=0$ là lúc vật bắt đầu rơi, chiều dương trùng với chiều chuyển động của vật.
\choiceTF[t]
{Thời gian từ khi vật được thả đến khi vật vừa chạm đất là $6\ s$}
{\True Tốc độ của vật ngay khi vừa chạm đất là $50 \ m/s$}
{\True Độ dịch chuyển của vật trong $2$ giây đầu là $20\ m$}
{\True Quãng đường của vật rơi trong giây thứ tư là $35\ m$}
\loigiai{
\begin{itemchoice}
\itemch Khi vật chạm đất thì quãng đường vật rơi bằng độ cao nơi thả vật $s=125 \ m$ nên thời gian rơi của vật từ khi thả đến khi vật vừa chạm đất là:
$$
t=\sqrt{\dfrac{2 \ s}{\ g}}=\sqrt{\dfrac{2.125}{10}}=5 \ s.
$$
\itemch 
Tốc độ của vật ngay khi vừa chạm đất là: $v=gt=10.5=50 \ m/s$.
\itemch 
Vì vật rơi tự do theo phương thẳng đứng, không đổi chiều nên độ dịch chuyển của vật bằng quãng đường vật rơi $d=s=\dfrac{1}{2} gt^2$

Trong thời gian $t=2 \ s$ đầu độ dịch chuyển của vật là: $d=\dfrac{1}{2} \cdot 10 \cdot 2^2=20 \ m$.
\itemch 
Quãng đường vật rơi trong 3 s đầu là: $s_3=\dfrac{1}{2} g \cdot t_3^2=\dfrac{1}{2} \cdot 10 \cdot 3^2=45 \ m$.
Quãng đường vật rơi trong 4 s đầu là: $s_{4}=\dfrac{1}{2} g \cdot t_{4}^2=\dfrac{1}{2} \cdot 10 \cdot 4^2=80 \ m$.
Quãng đường vật rơi trong giây thứ tư là: $\Delta s=s_{4}-s_3=80-45=35 \ m$.
\end{itemchoice}
}
\end{ex}
\begin{ex}
\immini[thm]{Thả một vật rơi tự do từ một đỉnh tháp cao $h$ (so với mặt đất) tại nơi có gia tốc trọng trường $g=10 \ m/s^2$. Tốc độ ngay khi vừa chạm đất của vật là $v=35 \ m/s$. Chọn mốc thời gian $t=0$ là lúc vật bắt đầu rơi, chiều dương trùng với chiều chuyển động của vật.}{\begin{tikzpicture}[scale=1,thick,>=stealth',draw=Mapcolor]

\draw[->] (0,0) -- (0,3.7);
\draw[->] (0,3) -- (2.7,3);

\draw (-0.1,1) -- (0.1,1);

\draw (-0.1,1.5) -- (0.1,1.5);
\draw (-0.1,2) -- (0.1,2);
\draw (-0.1,2.5) -- (0.1,2.5);

\draw (0.5,3.1) -- (0.5,2.9);

\draw (1,3.1) -- (1,2.9);
\draw (1.5,3.1) -- (1.5,2.9);
\draw (2,3.1) -- (2,2.9);
\draw[Mapcolor] (0,3) .. controls (1.5,2.5) and (2,1.5) .. (2.5,0.5);
\node at (0,4) {$s$};
\node at (3,3) {$t$};
\node at (-0.5,1) {$-20$};
\node at (-0.5,2) {$-10$};
\node at (0.5,3.5) {$1$};
\node at (-0.5,3) {$0$};
\draw[->] (5.5,0) -- (5.5,3.5);

\draw[->] (5.2,1) -- (8.5,1);

\draw (6,1.1) -- (6,0.9);
\draw (6.5,1.1) -- (6.5,0.9);
\draw (7,1.1) -- (7,0.9);
\draw (7.5,1.1) -- (7.5,0.9);
\draw (5.6,1.5) -- (5.4,1.5);
\draw (5.6,2) -- (5.4,2);
\draw (5.6,2.5) -- (5.4,2.5);
\draw (5.6,3) -- (5.4,3);

\node at (5.5,3.7) {$v$};
\node at (8.7,1) {$t$};
\node at (5.3,1.2) {$0$};
\draw[Mapcolor] (5.5,2) -- (7.5,2);
\node at (1.5,0) {Hình 1};
\node at (7,0) {Hình 2};
\end{tikzpicture}}
\choiceTF[t]
{Đồ thị quãng đường - thời gian của vật như hình 1}
{Đồ thị vận tốc - thời gian của vật như hình 2}
{\True Thời gian từ khi vật được thả đến khi vật vừa chạm đất là $3,5 \ s$}
{Quãng đường vật đi được sau $2\ s$ kể từ lúc bắt đầu rơi là $41,25 \ m$}
\loigiai{
\begin{itemchoice}
\itemch Quãng đường vật rơi phụ thuộc vào thời gian: $s=\dfrac{1}{2} g \cdot t^2=5 t^2$

Vì với mọi giá trị của $t \geq 0$ thì $s \geq 0$ nên đồ thị hàm số có dạng một nhánh của parabol qua gốc tọa độ và có bề lõm hướng lên.

$$
\text {itemch} \Rightarrow \text {Sai}
$$

Vận tốc của vật rơi tự do phụ thuộc vào thời gian: $v=g \cdot t=10 t \rightarrow$ Đồ thị hàm số có dạng một đường thẳng qua gốc tọa độ.
\itemch 
Khi vừa chạm đất vật có tốc độ $v=35 m/s$
Mặt khác tốc độ của vật tại thời điểm ${}^{t}$ kể từ $t=0$ là: $v=g \cdot t=10 t=35 \rightarrow t=3,5\ s$ 
\itemch 
Vật chạm đất sau thời gian $t=3,5 \ s$ nên độ cao nơi thả vật là: $h=5 t^2=5.3,5^2=61,25 \ m$ Quãng đường vật rơi sau 2 s là: $s_2=\dfrac{1}{2} g t_2^2=\dfrac{1}{2} \cdot 10 \cdot 2^2=20 \ m$.

\itemch
\end{itemchoice}
}
\end{ex}

\subsubsection*{PHẦN 3. Thí sinh trả lời từ câu 1 đến câu 6.}
\setcounter{ex}{0} 
\begin{ex}%Tailieuchuan
Hai vật ở độ cao $h_1$ và $h_2=10\ m$, được thả rơi tự do cùng một nơi. Thời gian rơi của vật thứ nhất bằng nửa thời gian rơi của vật thứ hai. Độ cao $h_1$ bằng bao nhiêu mét?
\shortans[oly]{2,5}
\loigiai{}
\end{ex}
\begin{ex}%Tailieuchuan
Trong $3\ s$ cuối cùng trước khi chạm đất, một vật rơi tự do đi được quãng đường bằng $\dfrac{1}{5}$ quãng đường toàn bộ mà nó rơi được. Lấy $g=10\ m/s^2$. Độ cao của vật lúc thả là bao nhiêu mét (làm tròn kết quả đến chữ số hàng đơn vị)? 
\shortans[oly]{4037}
\loigiai{}
\end{ex}
\begin{ex}%Tailieuchuan
Thả một hòn đá rơi từ miệng một cái hang sâu xuống đến đáy. Sau $3\ s$ kể từ lúc bắt đầu thả thì nghe tiếng hòn đá chạm vào đáy. Biết vận tốc truyền âm trong không khí là $340\ m/s$. Lấy $g=9,8\ m/s^2$. Chiều sâu của hang là bao nhiêu mét (làm tròn kết quả đến chữ số hàng phần mười)?
\shortans[oly]{40,7}
\loigiai{
}
\end{ex}

\begin{ex}%Tailieuchuan
Một vật rơi tự do khi chạm đất thì vật đạt vận tốc $40\ m/s$. Biết $g=10\ m/s^2$. Vật được thả rơi từ độ cao bao nhiêu mét? 
\shortans[oly]{80}
\loigiai{}
\end{ex}

\begin{ex}%Tailieuchuan
Một vật được thả rơi từ độ cao $1280\ m$ so với mặt đất. Lấy $g=10\ m/s^2$. Khi vận tốc của vật là $40\ m/s$ thì vật còn cách mặt đất bao nhiêu mét?
\shortans[oly]{1200}
\loigiai{}
\end{ex}
\begin{ex}%Tailieuchuan
Trong giây cuối cùng, một vật rơi tự do đi được một đoạn đường bằng nửa độ cao toàn phần $h$ mà vật rơi được kể từ vị trí ban đầu. Lấy $g=9,8\ m/s^2$. Độ cao $h$ bằng bao nhiêu mét? 
\shortans[oly]{57}
\loigiai{}
\end{ex}







